A integração de sensores inerciais (IMUs) com dados ópticos para maior precisão no controle de altitude e orientação trazem maior segurança e flexibilidade operacional das missões espaciais, a seguir será exibido mais uma tecnologia relacionada.

\section{Fusão de Fluxo Óptico e Medidas Inerciais}

Devido ao baixo consumo de energia, a utilização de câmeras como sensores para estimativa de movimento tem sido explorada cada vez mais, porém, a elevada quantidade de informações geradas exige um alto poder de processamento para estimar o deslocamento do sensor. Os métodos baseados exclusivamente em visão computacional foram desenvolvidos para rastrear a visão obtida pela câmera, mas existem limitações quando as características visuais do ambiente são inadequadas ou fora do que foi projetado pelo fabricante. Uma alternativa é a fusão de medições inerciais com os dados visuais, utilizando técnicas que combinam os dois tipos de informação para melhorar a estimativa de trajetória \cite{bloesch2014}.

Conforme explicado \cite{bloesch2014}, estudos, anteriores ao seu trabalho, desenvolveram abordagens para calibração conjunta entre sensores inerciais e câmeras, e destaca a importância da análise das características de observação do sistema. Os modelos que incorporam termos de erro visual em filtros de Kalman têm sido empregados para minimizar a incerteza na estimativa do movimento e os métodos como os baseados no filtro de Kalman não linear, utilizam medições inerciais para prever os dados do sensor, enquanto que as informações extraídas do fluxo óptico são inseridas como novos dados, permitindo uma atualização com maior precisão do que a iteração anterior.

Uma forma para reduzir a complexidade, consiste em evitar a inclusão direta das posições das imagens no estado do filtro, em vez disso, um termo de erro visual pode ser introduzido, permitindo uma estimativa eficiente sem haver a necessidade de armazenar múltiplas medições consecutivas do mesmo ponto. Dessa forma, este modelo facilita a implementação em tempo real, pois reduz a dimensionalidade do espaço de estados e minimiza o custo computacional. Além disso, um estudo de observabilidade não linear pode ser empregado para avaliar as propriedades do filtro e validar seu desempenho em dados experimentais \cite{bloesch2014}.
